\section{Conclusiones}

Este proyecto implementó un sistema completo de análisis de redes en C++17,
cubriendo siete métricas de centralidad aplicadas a dos redes reales con
características contrastantes: una red biológica dispersa y de gran escala,
y una red económica densa y pequeña.

Los experimentos confirman que \textbf{la representación importa}: la lista de
adyacencia permite analizar Yeast con 24 MB frente a los 340 MB de una matriz,
diferencia determinante para la factibilidad del análisis.

\textbf{Betweenness es el cuello de botella computacional}. Su complejidad
$O(|V|\cdot|E|)$ implica 109 segundos por ejecución en Yeast. El algoritmo de
Brandes es la única razón por la que este cálculo es viable; sin él tomaría días.
En contraste, PageRank y Eigenvector convergen en menos de 100 iteraciones,
siendo eficientes incluso en la red más grande.

El experimento de aristas revela una diferencia estructural fundamental:
\textbf{Yeast es robusta} (alta densidad, cambios mínimos ante perturbaciones),
mientras que \textbf{Trade es sensible en su periferia} (nodos poco conectados
amplifican el impacto de cada arista nueva). En términos aplicados, esto sugiere
que en el comercio internacional, integrar economías periféricas tiene mayor impacto
estructural que reforzar vínculos entre potencias ya conectadas.

Las métricas adicionales aportan dimensiones complementarias: Eigenvector Centrality
captura la \textit{influencia propagada} (estar conectado a nodos importantes) y
Clustering Coefficient mide la \textit{cohesión local} (formación de grupos o módulos).
Juntas con las cinco métricas base, ofrecen una visión multidimensional de la
importancia de cada nodo que ninguna métrica individual puede proporcionar por sí sola.
